%%%%%%%%%%%%%%%%%%%%%%%%%%%%%%%%%%%%%%%%%%%%%%%%%%%%%%%%%%%%%%%%%%%%%%
% Lớp 9 | Chương X. Một số hình khối trong thực tiễn | Hình trụ và hình nón
% Dạng 9H0?31-1: Diện tích xung quanh, thể tích hình trụ
%%%%%%%%%%%%%%%%%%%%%%%%%%%%%%%%%%%%%%%%%%%%%%%%%%%%%%%%%%%%%%%%%%%%%%


% nguồn: 2026-11-15_09_de_tham_khao_quan_3_de_1_va_de_2 (nhập 15/11/2026)
\begin{bt}%[9H0V31-1]%[ma:15b20df6]
% nguon: Đề tham khảo Tuyển sinh 10 TP.HCM 2025-2026 – Quận 3 – Đề 2 – Bài 4
Một hình chóp tứ giác đều có chiều dài cạnh đáy là $x\text{ (cm)}$ và chiều dài đường cao của mặt bên (trung đoạn) ngắn hơn cạnh đáy $2\text{ cm}$.
\begin{enumerate}[a)]
	\item Viết biểu thức $S_{xq}$ biểu diễn diện tích xung quanh của hình chóp theo $x$.
	\item Biết diện tích toàn phần của hình chóp tứ giác đều này là $207\text{ cm}^2$, hãy tính độ dài cạnh đáy $x$.
\end{enumerate}
\loigiai{
\begin{enumerate}[a)]
	\item Chiều dài đường cao của mặt bên (trung đoạn của hình chóp) là: $d = x - 2\text{ (cm)}$ (với $x > 2$).\\
	Hình chóp tứ giác đều có $4$ mặt bên là các tam giác cân có đáy $x$ và đường cao $d = x - 2$.\\
	Biểu thức biểu diễn diện tích xung quanh của hình chóp là:
	$S_{xq} = 4 \cdot \left(\dfrac{1}{2} x \cdot (x - 2)\right) = 2x(x - 2) = 2x^2 - 4x\text{ (cm}^2\text{)}$.
	\item Diện tích đáy của hình chóp tứ giác đều là: $S_{\text{đáy}} = x^2\text{ (cm}^2\text{)}$.\\
	Diện tích toàn phần của hình chóp là:
	$S_{tp} = S_{xq} + S_{\text{đáy}} = (2x^2 - 4x) + x^2 = 3x^2 - 4x\text{ (cm}^2\text{)}$.\\
	Theo đề bài, $S_{tp} = 207\text{ cm}^2$, ta có phương trình:
	$3x^2 - 4x = 207 \Leftrightarrow 3x^2 - 4x - 207 = 0$.\\
	Ta có $\Delta' = (-2)^2 - 3\cdot(-207) = 4 + 621 = 625 = 25^2 > 0$.\\
	Vì $x > 2$ nên phương trình có nghiệm hợp lệ duy nhất:
	$x = \dfrac{2 + 25}{3} = \dfrac{27}{3} = 9\text{ (cm)}$.\\
	Vậy độ dài cạnh đáy của hình chóp tứ giác đều là $9\text{ cm}$.
\end{enumerate}
}
\end{bt}
